\src{188}
\begin{equation}\tag{7.2}
  \symup{f^{*}:H^{*}(X)\longrightarrow H^{*}(Y)}
\end{equation}
Comme il s'agit d'algèbres de Poincaré, on obtient par transposition un
homomorphisme d'espaces vectoriels :
\begin{equation}\tag{7.3}
  \symup{f_{*}:H^{*}(Y)\longrightarrow H^{*}(X)}
\end{equation}
On a vu au n\textsuperscript{o} 4 que l'effet de $\symup{f^{*}}$ sur les classes de
cohomologie correspondant à des cycles non singuliers s'interprète
géométriquement en prenant les classes de cohomologie qui correspondent à leurs
images inverses. Il importe dans le cas actuel de montrer que (7.3) correspond
de même à l'opération \fquote{image directe} de cycles. Cela résulte, sous des
conditions de non singularité convenables du moins, du cas particulier suivant :

\result{THÉORÈME}{4} \uline{Si} $\symup{f}$ \uline{est l'application identique d'une sous-variété non
singulière} $\symup{Y^{m}}$ \uline{de} $\symup{X^{n}}$ \uline{dans} $\symup{X^{n}}$, \uline{alors, désignant par} $\symup{1_Y}$
\uline{l'élément unité de} $\symup{H(Y)}$ \uline{on a}
\begin{equation}\tag{7.4}
  \symup{f_{*}(1_Y)=P_X(Y)}
\end{equation}
\uline{où le deuxième membre est la classe de cohomologie dans} $\symup{X}$ \uline{associée à} $\symup{Y}$.

Cette formule est équivalente à
\begin{equation}\tag{7.4 bis}
  \symup{\bigl\langle\xi^{m,m},P_X(Y)\bigr\rangle\varepsilon_Y
  =f^{*}(\xi^{m,m})
  \qquad
  \bigl(\xi^{m,m}\in H^{m}(X,\underline{\Omega}_X^{m})\bigr)}
\end{equation}
\footnote{Le texte imprimé porte $\symup{f_{*}(\xi^{m,m})}$. Puisque
$\symup{\xi^{m,m}\in H^{m}(X,\underline{\Omega}_X^{m})}$ et que les deux membres appartiennent à
$\symup{H^{m}(Y,\underline{\Omega}_Y^{m})}$, l'application requise est $\symup{f^{*}}$.} où
$\symup{\varepsilon_Y}$ est l'élément fondamental dans $\symup{H^{m}(Y,\underline{\Omega}_Y^{m})}$, et
constitue, dans le cas des variétés projectives non singulières, une nouvelle
définition de la classe de cohomologie associée à $\symup{Y}$. Pour démontrer le
théorème 4, on interprète grâce au théorème 3 la transposée de l'homomorphisme
\[
  \symup{H^{m}(X,\underline{\Omega}_X^{m})
  \longrightarrow H^{m}(Y,\underline{\Omega}_Y^{m})
  =H^{m}(X,\underline{\Omega}_Y^{m})}
\]
comme l'homomorphisme
\begin{equation}\tag{7.5}
\symup{H^{n-m}(X,\underline{\Omega}_X^{n-m})
  \xrightarrow{\sim}
  \operatorname{Ext}_{\Osh_X}^{n-m}(X;\underline{\Omega}_X^{m},\underline{\Omega}_X^{n})
  \longleftarrow
  \operatorname{Ext}_{\Osh_X}^{n-m}(X;\underline{\Omega}_Y^{m},\underline{\Omega}_X^{n})
  \xrightarrow{\sim}
  \operatorname{Hom}_{\Osh_X}(X;\underline{\Omega}_Y^{m},\underline{\Omega}_Y^{m})}
\end{equation}
On vérifie que l'élément $\symup{1_Y}$ du dual de $\symup{H^{m}(Y,\underline{\Omega}_Y^{m})}$ s'identifie
à l'élément du membre de droite correspondant à l'endomorphisme identique de
$\symup{\underline{\Omega}_Y^{m}}$, et d'autre part que l'image de cet élément dans
$\symup{H^{n-m}(X,\underline{\Omega}_X^{n-m})}$ est bien $\symup{P_X(Y)}$. Ces compatibilités, qui auraient
pu être données dès le n\textsuperscript{o} 4, peuvent s'énoncer, et sont
valables en fait, pour des variétés non singulières quelconques, la deuxième

